\documentclass[10pt]{article}

\usepackage[a4paper, total={7in, 10.5in}, includehead]{geometry}
\usepackage{parth-article}
\usepackage{sensible-macros}
\usepackage[fixed]{fontawesome5}
\addbibresource{~/Notebook/.meta/zotero-library.bib}

\setmainfont{baskervaldx}[Numbers={Proportional, OldStyle}]

\begin{document}
\title{
	\vspace{-1em}Parth Shimpi\normalsize\orcidlink{0009-0008-2066-7474}
	\raisebox{-0.1em}{\huge\(\vert\)}
	\normalsize\ \raisebox{0.1em}{List of publications}
}
\date{}
\maketitle
\thispagestyle{fancy}
\begin{description}
	\item[``Torsion pairs and 3-fold flops"] (2025).

	      \emph{E-print available at}
	      \href{https://arxiv.org/abs/2502.05146}{arXiv:2502.05146}.

	      The paper classifies t-structures in the derived category of a 3-fold flopping
	      contraction that have small cohomological spread with respect to the standard
	      algebraic heart. This is a foundational step towards describing all
	      t-structures and spherical objects in the derived category. Immediate
	      applications of the results, to the classification of spherical modules,
	      bricks, and semibricks are given.

	      In equivalent algebraic terms, the paper describes the complete lattice of
	      torsion classes for an affine preprojective algebra. For this, new
	      convex-geometric techniques are introduced which enable an arbitrary
	      t-structure to be approximated by algebraic ones. The techniques have
	      applications to other classes of finite-dimensional algebras.

	\item[``Simple complexes on a flopped curve''] (2025).

	      \emph{E-print available at}
	      \href{https://arxiv.org/abs/2508.06437}{arXiv:2508.06437}.

	      The paper classifies, completely and without any caveats, all spherical objects
	      and stability conditions on a single-curve flopping contraction. Consequently,
	      long-standing conjectures about connectedness and contractibility of the
	      stability manifold are shown to have affirmative answers in this context. In
	      fact a stronger result is established, which classifies all objects satisfying
	      a natural cohomological constraint. Algebraically, this is the first
	      classification of tilting complexes on a g-affine algebra. This is proof of
	      concept that the results of \emph{``Torsion pairs and 3-fold flops''}, with
	      appropriate geometric insight, can extend to global classification results.

	\item[``Kleinian orbifolds, cohomological Hall algebras, and Yangians'']
	      (2025),

	      \emph{joint with} Francesco Sala \emph{and} Olivier Schiffmann.
	      \emph{In preparation}.

	      The paper computes cohomological Hall algebras for every Abelian category
	      appearing as the heart of a stability condition on the minimal resolution of a
	      Kleinian singularity. In particular CoHAs are computed for a new class of
	      surfaces with non-commutative structure, namely those obtained as orbifold
	      resolutions. The key technical tool is the approximation-based analysis of
	      ``Torsion pairs and 3-fold flops'', which allows for the natural t-structure on
	      these surfaces to be approximated by well-understood algebraic ones.

	      Preprojective algebras and minimal surfaces are respectively the standard
	      algebraic and geometric resolutions of a Kleinian singularity, and their
	      cohomological Hall algebras are known to coincide with distinct `positive
	      halves' of the Maulik--Okounkov Yangian. Orbifold resolutions form a family of
	      part-algebraic, part-geometric resolutions that interpolate between the two
	      extremes. The computation of their cohomological Hall algebras recovers new
	      halves of the Yangian.

\end{description}
\end{document}
